\par\bigskip
\begingroup\small
\begin{center}BEMERKUNGEN ZU DEN ABSCHNITTEN V, VI, VII UND ZUM BRIEFWECHSEL.\end{center}
\begin{center}\textls[120]{Die Abschnitte} [V.] \textls[120]{und} [VI.], S. 400—419.\end{center}
\indent Im Nachlaß von \textsc{Gauss} finden sich \textls[120]{fünf} verschiedene Fassungen für den Anfang einer Abhandlung »\textit{Convergenz der Reihen, in welche die periodischen Functionen einer veränderlichen Grösse entwickelt werden}«.

\indent Die vier ältern Fassungen haben das gemeinsam, daß sie mit der Erörterung des Begriffs der »\textls[120]{Konvergenz einer Reihe}« beginnen und über diese Erörterung auch nicht wesentlich hinauskommen\textsuperscript{*}). Wir haben im Abschnitt [V.] die beiden ersten Artikel des \textls[120]{vierten} und ausführlichsten dieser ältern Entwürfe wiedergegeben; in der Handschrift folgen noch zwei Artikel, die Erörterungen über die Einteilung der Funktionen in explizite und implizite enthalten. Der Abdruck dieser beiden Artikel 3. und 4. erschien zwecklos. Dagegen haben wir die am Schluß des art. 1., oben S. 403, erwähnte Erweiterung, die sich auf einem einzelnen Blatte gefunden hat, als \textls[120]{Anhang} abgedruckt. Man wird in den Erörterungen dieses art. 1. und des Anhangs die Anfänge der Betrachtungen wiedererkennen, die \textsc{Paul du Bois-Reymond} als \textls[120]{Infinitärkalkül}\linebreak bezeichnet hat\textsuperscript{**}). Bei der Einführung der iterierten Logarithmen neben der Potenz, wodurch eine »nähere Anschmiegung erreicht werden kann« (oben S. 403), im \textls[120]{Anhang}, hätte \textsc{Gauss} auf \textsc{Abel}\textsuperscript{***}) verweisen können, falls er auch diese Untersuchungen nach 1831 angestellt haben sollte, was für die artt. 1. und 2. durch die Erwähnung der Anzeige von 1831 im art. 2., oben S. 405, feststeht. Während bei \textsc{Gauss} nur die einfache Exponentialgröße \(e^{\alpha n}\) auftritt, hat \textsc{P. du Bois-Reymond} a. a. O. seine Skalen noch durch iterierte Exponentialfunktionen erweitert. Wir bemerken auch, daß \textsc{P. du Bois-Reymond} ausdrücklich auf die Wichtigkeit des Infinitärkalküls für die Theorie der \textsc{Fourier}schen Reihen hinweist; es
\par\medskip\noindent\rule{.24\linewidth}{.4pt}\par\smallskip
\noindent\hspace*{\parindent}\textsuperscript{*}) Der erste Entwurf enthält 3 Quartseiten, der zweite 3 Oktavseiten, der dritte 6 Quartseiten, der vierte 13 Quartseiten.
\par\noindent\hspace*{\parindent}\textsuperscript{**}) Siehe Annali di Matematica, 2. Serie \textbf{4}, 1870, S. 339; \textsc{Crelles} Journal für Mathematik \textbf{74}, 1872, S. 294; Mathem. Annalen \textbf{8}, 1875, S. 363; \textbf{11}, 1877, S. 149; vergl. auch A. \textsc{Pringsheim}, Mathem. Annalen \textbf{35}, 1890, S. 302 und \textit{Vorlesungen über Zahlen- und Funktionenlehre} I, 1, 1916, S. 224 ff.
\par\noindent\hspace*{\parindent}\textsuperscript{***}) \textsc{Crelles} Journal für Mathematik \textbf{3}, 1828, S. 79, Oeuvres de N. H. \textsc{Abel}, nouvelle édition, I, 1881, S. 399.
\clearpage\SourceLeaf{445}\ArithmeticHead{438}{L}{ANALYSIS. NACHLASS.}
\noindent zeigt dies, daß die Untersuchungen des art. 1. und des \textls[120]{Anhangs} in naher Beziehung zu dem in dem Titel der Abhandlung bezeichneten Gegenstande stehen.

\indent Die Erörterungen des art. 2. des Abschnitts [V.] werden in der \textls[120]{fünften} Fassung, die wir in den artt. 1.—8. des Abschnitts [VI.] unverkürzt wiedergeben, weiter ausgebaut. Im Eingang zum art. 1. dieses Abschnitts wird von der in der Überschrift bezeichneten Aufgabe gesagt, daß sie »recht eigentlich« in die Analysis der komplexen Größen gehöre; es wird dann in den artt. 2.—5. in Auseinandersetzungen eingetreten, für die die Worte gelten, mit denen \textsc{Gauss} den art. 5. seiner \textit{Jubiläumsschrift} (Werke III, S. 79) abschließt, nämlich, daß ihr »Gegenstand die nach der Stetigkeit zusammenhängenden Größenkombinationen sind«, die zwar »einem höhern von Räumlichem unabhängigen Gebiete der allgemeinen abstrakten Größenlehre« angehören, »in welchem man sich \({.}\;{.}\;{.}\) [aber] nicht bewegen kann, ohne eine von räumlichen Bildern entlehnte Sprache«. Die in den artt. 6.—8. enthaltenen Untersuchungen über Integrale zeigen, daß \textsc{Gauss} beabsichtigt hat, die Untersuchungen der artt. 2.—5. auf die Lehre von der Integration im komplexen Gebiete anzuwenden\textsuperscript{*}). Leider hat sich im Nachlaß keine Aufzeichnung vorgefunden, die sich auf diese Anwendung bezieht, und auch keine, in der von der weitern Anwendung der Analysis der komplexen Größen auf den in der Überschrift bezeichneten Gegenstand die Rede wäre. Die Aufzeichnung, die wir dem Abschnitt [VI.] als art. 9. hinzugefügt haben, bezieht sich zwar auf einen hierher gehörigen Gegenstand, sie läßt aber auch kaum erkennen, was \textsc{Gauss} beabsichtigt haben mag. Dagegen werden wir nachher an der Hand der Stücke [2.]—[4.] des \textit{Briefwechsels} zeigen, daß die Wurzel der geplanten Abhandlung, von der uns in den Abschnitten [V.] und [VI.] nur einige einleitende Kapitel überkommen sind, in den im Abschnitt [VII.] zusammengestellten Aufzeichnungen enthalten ist, zu deren Deutung und Erklärung wir uns jetzt wenden.
\par\medskip\noindent\rule{.24\linewidth}{.4pt}\par\smallskip
\noindent\hspace*{\parindent}\textsuperscript{*}) R. \textsc{Dedekind} berichtet in seiner Lebensbeschreibung \textsc{Riemanns} (siehe \textsc{Riemanns} Werke, 2. Auflage, 1891, S. 545), daß \textsc{Gauss}, als \textsc{Riemann} ihn 1851 vor der mündlichen Doktorprüfung besuchte, geäußert habe, er bereite seit Jahren eine Schrift vor, die denselben Gegenstand behandele, wie \textsc{Riemanns} Inauguraldissertation, sich aber freilich nicht darauf beschränke. Offenbar hat \textsc{Gauss} dabei an die hier in Rede stehende Abhandlung »\textit{Bestimmung der Convergenz der Reihen u.s.w.}« gedacht.
\par\endgroup
\ArithmeticHead{438}{L}{ANALYSIS. NACHLASS.}
